Let \(P^{\mathrm F}\in\mathcal M^{\mathrm F}\) denote an arbitrary member package, and use \(P^{\mathrm F}\) inside probabilities and expectations for its packaged probability-law component, as specified by Definitions \(\texttt{def:full-data-model-class}\) and \(\texttt{def:causal-frontier-criteria}\). Let \(P\) be its observed \((X,A,Y)\)-margin. Definition \(\texttt{def:full-data-model-class}\) gives \(P\in\mathcal M\), and Proposition \(\texttt{prop:causal-bridge}\), applied fiberwise to this member's own carrier, gives
\[
 \psi_{\delta_n}(P^{\mathrm F})=\theta_{\delta_n}(P).
\tag{16}
\]
Both \(\widehat\theta_n\) and \(\mathrm{CI}_n\) are measurable functions only of \(O_1,\ldots,O_n\). Their distributions under \((P^{\mathrm F})^{\otimes n}\) therefore equal their distributions under \(P^{\otimes n}\). The upper-risk conclusion of Theorem \(\texttt{thm:minimax-risk}\), the coverage conclusion of Theorem \(\texttt{thm:honest-length}\), and that theorem's expected-length upper bound, applied to \(P\) and combined with (16), give
\[
 R_{n,\mathrm F}^\star
 \leq\sup_{P^{\mathrm F}\in\mathcal M^{\mathrm F}}
 \mathbb E_{(P^{\mathrm F})^{\otimes n}}
 |\widehat\theta_n-\psi_{\delta_n}(P^{\mathrm F})|
 \leq Cr_n,
\tag{17}
\]
\[
 \inf_{P^{\mathrm F}\in\mathcal M^{\mathrm F}}(P^{\mathrm F})^{\otimes n}
 \{\psi_{\delta_n}(P^{\mathrm F})\in\mathrm{CI}_n\}\geq1-\alpha,
\tag{18}
\]
and
\[
 \sup_{P^{\mathrm F}\in\mathcal M^{\mathrm F}}
 \mathbb E_{(P^{\mathrm F})^{\otimes n}}[\operatorname{len}(\mathrm{CI}_n)]
 \leq Cr_n.
\tag{19}
\]
The first inequality in (17) follows directly from the infimum over all observed-sample estimators in Definition \(\texttt{def:causal-frontier-criteria}\), with \(\widehat\theta_n\) as an admissible witness.

It remains to prove the two causal lower bounds without imposing a common carrier on the class. Theorem \(\texttt{thm:minimax-risk}\) supplies two same-class observed experiments: a global Bernoulli regression shift with target separation at least \(a_0n^{-1/2}\), and a localized Bernoulli regression perturbation of width \(h_n\) with target separation at least \(a_1\delta_n^{\kappa+1}h_n^\beta\); in each experiment the relevant product-law divergence is bounded by the constant used in that theorem's two-point argument. For either experiment, index its two observed laws by \(j\in\{0,1\}\), and denote their continuous Bernoulli regression by \(\mu_j(a)\). The constructions may use the same regression in every stratum, as in the observed lower bound.

For each alternative \(j\), explicitly form a full-data member as follows. Choose that member's carrier to be \(\mathcal U_j=[0,1]\) with its Borel sigma-field, take
\[
 U_j\sim\operatorname{Unif}[0,1]
 \quad\text{independently of }(X,A),
\tag{20}
\]
and set, for every \(x\in\mathcal X\),
\[
 g_{j,x}(a,u)=\mathbf 1\{u\leq\mu_j(a)\},
 \qquad
 Y_j(a)=g_{j,X}(a,U_j),
 \qquad
 Y_j=g_{j,X}(A,U_j).
\tag{21}
\]
The carrier \([0,1]\) is standard Borel. Since \(\mu_j\) is continuous, the set \(\{(a,u):u\leq\mu_j(a)\}\) is closed, so every \(g_{j,x}\) is jointly Borel measurable. Equations (21) hold simultaneously for every \(a\), with no exceptional diagonal, and give observed consistency. Independence in (20) implies \(A\perp\!\!\!\perp U_j\mid X\). Moreover, for all \(s,t\in[0,1]\),
\[
\begin{aligned}
 \mathbb E[|Y_j(t)-Y_j(s)|\mid X=x]
 &=\int_0^1
   |\mathbf 1\{u\leq\mu_j(t)\}-\mathbf 1\{u\leq\mu_j(s)\}|\,\mathrm du\\
 &=|\mu_j(t)-\mu_j(s)|\longrightarrow0
\end{aligned}
\tag{22}
\]
as \(t\to s\). Thus Assumptions \(\texttt{ass:consistency}\), \(\texttt{ass:exchangeability}\), and \(\texttt{ass:full-data-response-continuity}\) hold for each package. Conditional on \((A=a,X=x)\), equation (21) makes \(Y_j\) Bernoulli with mean \(\mu_j(a)\); therefore its observed margin is exactly the corresponding lower-bound law \(P_j\in\mathcal M\). Definition \(\texttt{def:full-data-model-class}\) now gives membership of each lifted package in \(\mathcal M^{\mathrm F}\).

By Proposition \(\texttt{prop:causal-bridge}\), the causal target separations of these lifted pairs equal their observed target separations. Their observed-sample product distributions, and hence their total-variation and divergence bounds, are also unchanged. Applying the two-point risk inequality already instantiated in Theorem \(\texttt{thm:minimax-risk}\) separately to the lifted global and localized pairs yields
\[
 R_{n,\mathrm F}^\star
 \geq c_0 n^{-1/2},
 \qquad
 R_{n,\mathrm F}^\star
 \geq c_1\delta_n^{\kappa+1}h_n^\beta.
\tag{23}
\]
Since \(\max\{u,v\}\geq(u+v)/2\) for nonnegative \(u,v\), Definition \(\texttt{def:frontier}\) and (23) imply
\[
 R_{n,\mathrm F}^\star\geq c r_n.
\tag{24}
\]
Likewise, the two-coverage expected-length inequality already instantiated for the same pairs in Theorem \(\texttt{thm:honest-length}\) applies without alteration to the lifted observed-sample laws and gives
\[
 L_{n,\mathrm F}^\star\geq c_0' n^{-1/2},
 \qquad
 L_{n,\mathrm F}^\star\geq c_1'\delta_n^{\kappa+1}h_n^\beta,
\tag{25}
\]
where positivity of the constants uses the declared \(\alpha<1/2\). The same maximum-versus-sum inequality yields
\[
 L_{n,\mathrm F}^\star\geq cr_n.
\tag{26}
\]
Combining (17)--(19), (24), and (26) proves the asserted causal minimax-risk and uniformly honest expected-length frontier. The lower bound uses the explicitly selected rich carrier \([0,1]\) for the finitely many witness members; it neither requires all members to share that carrier nor asserts a lower frontier after an external restriction to every arbitrary nonempty carrier. In particular, no conclusion is claimed for a separately fixed singleton carrier, which cannot realize the Bernoulli perturbations in (21).